\scanpage{11}
conique $x^2-y^2=0$, mais celle-ci a l'origine comme pt singulier. C'est une réunion des deux droites diagonales.

\margin{bref !}
On voit \uline{donc} que si ${\unsure{\zeta|\xi}}^{4}-1$ est inv. i.e. \deleted{\unsure{$\zeta$|$4$}}$(\alpha^2-4)$ inv., alors \deleted{il y a ex} une forme quadratique sur $E$ quasi-invariante par $\Phi$ \uline{et non dég.} sur chaque fibre est en fait invariante par $\Phi$, et c'est une combinaison linéaire des formes $z^2$ et $xy$ (avec les coordonnées choisies\unsure{\}|)}, \deleted{qui} où on a dû d'abord faire un changt de base quadratique vers
\[
S'=\operatorname{Spec}\underline O_S[T]/(T^2-\alpha T+1)\quad\ldots\,)
\]
\[
Q(x,y,z)=az^2+bxy\tag*{(5)}
\]
de \deleted{déterminant} matrice
\[
\begin{pmatrix}0&b&0\\b&0&0\\0&0&2a\end{pmatrix}
\]
de déterminant
\[
\delta(Q)=2ab^2\tag*{(6)}
\]
donc de déterminant divisé
\[
\delta'(Q)=ab^2\tag*{(7)}
\]
\unsure{(qui|Qui} est \deleted{\unsure{inversible}}\inserted{inv.} \unsure{sss|ssi} $a,b$ le sont. La condition que $Q$ s'annule sur les sommets de $P$ s'écrit alors
\[
a+b=0\tag*{(8)}
\]
donc \deleted{\unsure{les}} $\exists!$ \deleted{\unsure{f.}} \unsure{côniques|cônique} circonscrites à $P$ invariantes sous $\mathbb D_\infty$, \deleted{sont} d'équation
\[
z^2-xy=0\tag*{(9)}
\]
Ou encore, en coordonnées affines, $\exists!$
% archive sheet number: 10
